\documentclass{article}
\usepackage{amsmath}
\title{A recovery certificate can consume the capacity it certifies}
\author{}
\date{}
\begin{document}
\maketitle
\paragraph{Scope.}
Held-in development case for I1 and the L2 evidence/action boundary.
The parent replay establishes that a source-qualified removal before a
relay's first send permits the recovery controller to distinguish a lost
reply from one already charged against the receiver. Here that evidence
travels through the actual selected SubT broker and radio code, competing
with two existing report duties. All plants are simulations. This is
not a field protocol, independent mission, or architecture disposition.

\paragraph{Source and interface.}
The same pinned scorer as the parent produces two exact, distinct
keys. Both are entered and read in a SQLite archive; three native bound
callbacks stage their handles with two holders $P,Q$ of $A$ and one holder
$R$ of $B$. Another score closes the scorer and independent traffic
still crosses the broker. At tick~4, $P$ is removed before its first
possible attempt or after its actual deposit. The witness emits either
\texttt{E} (the selected removal preceded that attempt), or \texttt{X}
(the relay is now absent, without the event-order guarantee). The
post-attempt receiver notification for $A$ can be hidden. A silent
pre-attempt removal yields neither token. The committed native ingress
effect of $P$ is not undone by its removal.

The fixture generates and binds the one-byte tokens to this exact
request and observed event, uses a trusted named endpoint, and has no
old tokens, forged sender, or unmodeled traffic. This does not establish
that a stock deployment offers such a one-byte signed message. The
alternative local channel provides exactly the same facts without a
radio debit, an explicit resource assumption, not an implementation of a
local sensor. For the main sweep, transport, configured caps, and the
scheduled receive tick for the control report are known at the cut.
The checker reads only the current qualified receipt notice, the
qualified removal, and the weaker absence report. It does not read
the native receive counter, the scorer truth as a current receipt,
or the hidden world identifier. A second sweep hides the receive
tick and retains all compatible prefix rows.

\paragraph{Exact cut calculation.}
The inherited charges are per sender and receiver, with the same sender
radio limit used for both tests. $Q$ sends $2$ bytes at cap~$4$, and
$R$ sends $3$ at cap~$5+w$ ($w=0$ or $1$). The native one-second
window is four ticks and retains debits at equality. For a real
\texttt{E} received at tick $s$, let $c_s(t)=[t\le s+4]$ at a
subsequent send tick $t$ (or zero off-radio). At the recovery cut, the
exact complete service test for sending $Q$ at $u$ and $R$ at $v$,
$u\le v\le u+4$, is
\[
       c_s(u)+2\le4\quad\land\quad c_s(v)+2+3\le5+w.
\]
The first inequality always holds in the declared envelope. The second
is equivalent to $v>s+4$ or $w\ge1$ for the charged certificate.
In reverse order the debit from $B$ still lives when $A$ sends
(for all tested send times), and $3+2>4$ regardless of the certificate.
No schedule within this horizon can improve on sending each required
copy once: erase unsuccessful, unnecessary, or repeated attempts in
a fixed history and the load at future attempts cannot increase. In the
native fixture there are no new observations offered between the two
send slots, so the same deletion applies to all histories in the
qualified \texttt{E} cell. The Lean file checks the arithmetic, not the
authenticity of the byte or the binding to the upstream implementation.

\paragraph{Discriminating cases.}
With due tick~8 and $w=0$, local evidence permits both duties at tick~8;
the same evidence sent over the radio at tick~4 does not permit both:
the old debit is still live. If $A$ is due at~8 and $B$ may wait until~9,
send $Q$ at~8 and $R$ at~9. No other tested sequence meets those
deadlines. Moving the same report's arrival from~4 to~5 removes that
option, and pooling ticks~4 and~5 without a reliable timestamp
requires declining that option even in the actual tick~4 history.
Raising $R$'s cap by one restores $Q,R$ even with the late charged
certificate; it does not restore the reverse order. In the hidden
post-deposit history, sending the weaker one-byte \texttt{X} cannot
certify absence of a prior charge. It can even cause $B$ alone to
miss at tick~8: $2+1+3>5$. Source removal, current membership,
timestamp, and evidence cost are therefore separately relevant.

\paragraph{Verification and kernel.}
The unchanged shared checker selects the native continuation. All
candidate schedules and the actual selected continuation are tested
against the native scorer/broker and a direct debit replay. Full
service and protecting only $B$ are checked separately; if neither
is admitted, the harness' default attempted send is logged as
unproved, not counted as service. The same-information direct method
ties. No result here requires a full packet-log replay at the cut, but
the channel, scope, and relevant expiry cannot be forgotten. Freshness,
effect termination, authorization to change a deadline, more general
sensing/route costs, and a true independently chosen mission remain
open. The replay does not silently allocate a separate receive queue
or assume the certificate's transmission is free.
\end{document}
